\section{Kth Largest Element in a Stream}
\label{sec:heap-kth-largest-stream}

\problemstatement{Design a class initialized with an integer $k$ and an
initial array $\textit{nums}$, supporting $\textit{add}(\textit{val})$:
add $\textit{val}$ to the running stream and return the $k$-th largest
element among \textbf{all} elements seen so far (including the initial
array and every value added).}

\begin{patternbox}
This is Section~\ref{sec:heap-kth-largest}'s ``Kth Largest Element in an
Array'' made \textbf{incremental}: instead of computing the answer once
for a fixed array, we must maintain it continuously as new elements
stream in, one at a time. The keyword shift from a single static query to
a class with a repeated $\textit{add}$ method is the signal that the
size-$k$ min-heap should be built \emph{once} and \emph{kept alive}
across calls, rather than recomputed from scratch on every query.
\end{patternbox}

\subsection*{Understanding the problem carefully}

The min-heap-of-size-$k$ technique from
Section~\ref{sec:heap-kth-largest} already processes elements one at a
time and remains correct at \emph{every intermediate point} during that
processing -- the heap, after processing any prefix of elements, always
holds the $k$ largest elements seen \emph{so far}, with the $k$-th largest
sitting at the root. This means no new derivation is needed: we simply
keep the heap object alive as part of the class's state, seed it with the
initial array, and apply the exact same push-then-conditionally-pop step
to each newly added value.

\begin{ideabox}[Reuse: Make the Heap Persistent]
On construction: build a min-heap from the initial $\textit{nums}$,
trimming it down to size $k$ using the same push-then-pop-if-oversized
rule as Algorithm~\ref{sec:heap-kth-largest}. On each $\textit{add}(val)$:
apply that same rule to $\textit{val}$ alone, then return the heap's
current root.
\end{ideabox}

\subsection*{Algorithm}

\begin{enumerate}
  \item \textsc{Constructor}($k$, \textit{nums}): for each $x$ in
        \textit{nums}, push $x$ onto the min-heap; if size exceeds $k$,
        pop the minimum.
  \item \textsc{Add}($\textit{val}$): push $\textit{val}$ onto the
        min-heap; if size exceeds $k$, pop the minimum; return the
        heap's current top.
\end{enumerate}

\begin{algorithm}[H]
\caption{Kth Largest Element in a Stream}
\begin{algorithmic}[1]
\Procedure{Constructor}{$k$, \textit{nums}}
  \State $\textit{heap} \gets$ empty min-heap
  \For{each $x$ in \textit{nums}}
    \State push $x$ onto \textit{heap}
    \If{$|\textit{heap}| > k$} pop the minimum \EndIf
  \EndFor
\EndProcedure
\Procedure{Add}{\textit{val}}
  \State push \textit{val} onto \textit{heap}
  \If{$|\textit{heap}| > k$} pop the minimum \EndIf
  \State \Return top of \textit{heap}
\EndProcedure
\end{algorithmic}
\end{algorithm}

\subsection*{Dry run}

\dryrun{Take $k=3$, initial $\textit{nums}=[4,5,8,2]$, then call
$\textit{add}(3)$, $\textit{add}(5)$, $\textit{add}(10)$.}

\begin{itemize}
  \item Construct: push $4,5,8,2$ one at a time, trimming to size $3$
        each time it exceeds $3$. After all four: heap holds the $3$
        largest of $\{4,5,8,2\}$, which are $\{4,5,8\}$ (root $=4$).
  \item $\textit{add}(3)$: push $3$; heap $=\{3,4,5,8\}$, size $4>3$:
        pop min ($3$). Heap $=\{4,5,8\}$. Return $4$.
  \item $\textit{add}(5)$: push $5$; heap $=\{4,5,5,8\}$, size $4>3$:
        pop min ($4$). Heap $=\{5,5,8\}$. Return $5$.
  \item $\textit{add}(10)$: push $10$; heap $=\{5,5,8,10\}$, size $4>3$:
        pop min ($5$). Heap $=\{5,8,10\}$. Return $5$.
\end{itemize}

\subsection*{C++ implementation}

\begin{lstlisting}
#include <bits/stdc++.h>
using namespace std;

class KthLargest {
    int k;
    priority_queue<int, vector<int>, greater<int>> minHeap;

public:
    KthLargest(int k, vector<int>& nums) : k(k) {
        for (int x : nums) {
            minHeap.push(x);
            if ((int)minHeap.size() > k) minHeap.pop();
        }
    }

    int add(int val) {
        minHeap.push(val);
        if ((int)minHeap.size() > k) minHeap.pop();
        return minHeap.top();
    }
};

int main() {
    vector<int> nums = {4, 5, 8, 2};
    KthLargest kl(3, nums);

    cout << kl.add(3) << endl;
    cout << kl.add(5) << endl;
    cout << kl.add(10) << endl;
    return 0;
}
\end{lstlisting}

\begin{complexitybox}
\textbf{Time Complexity.} Construction costs $O(n \log k)$ for the
initial $n$ elements. Each subsequent $\textit{add}$ call costs
$O(\log k)$.
\textbf{Space Complexity.} $O(k)$, since the heap never exceeds size $k$.
\end{complexitybox}

\begin{takeawaybox}
Many ``streaming'' versions of array problems require no new algorithmic
idea at all -- only recognizing that the data structure from the static
version (here, a size-$k$ min-heap) is already incremental by
construction, and simply needs to be kept alive as persistent state
rather than rebuilt on every query.
\end{takeawaybox}
